\section{Spatial Witnesses and the Limits of Pooling}
\label{sec:spatial}
The preceding construction is valid for any chosen summaries. Its
interpretation nevertheless depends on how they were formed. This
section isolates the spatial issue and gives an exact representation
when max pooling is too coarse.

\subsection{Separate-site and same-site evidence}
\label{sub:spatial-witnesses}
For a query $u\in\U$, define two image retrieval sets:
\begin{equation}\label{eq:spatial-extents}
  \begin{aligned}
    G_{\mathrm{img}}(u)
         & =\{x\mid \forall j\ \exists p\in P_{x}: u_{j}\leq z_{x,p,j}\}, \\
    H(u) & =\{x\mid \exists p\in P_{x}\ \forall j: u_{j}\leq z_{x,p,j}\}.
  \end{aligned}
\end{equation}
The first is $G_{I}(u)$ for max-pooled summaries. The second requires a
single site to witness every coordinate. It is also the image projection
of the extent in the site-level context.

\begin{proposition}[Spatial relaxation under max pooling]
  \label{prop:spatial-relaxation}
  For every query, $H(u)\subseteq G_{\mathrm{img}}(u)$. Equality holds for
  queries with at most one positive coordinate, but can fail for two
  positive coordinates. In general, $H$ need not take joins to intersections.
\end{proposition}
\begin{proof}
  A site witnessing every coordinate supplies a witness for each coordinate
  separately. With one positive coordinate, a site attaining its maximum
  witnesses that coordinate, while nonnegativity satisfies the others.
  For a zero query, every image qualifies since each $P_{x}$ is nonempty.
  For strictness, take one image with codes $(2,0)$ and $(0,2)$. It belongs
  to $G_{\mathrm{img}}((2,2))$ but not to $H((2,2))$. It also belongs to
  both $H((2,0))$ and $H((0,2))$, proving failure of join-to-intersection
  preservation.
\end{proof}

The example prevents treating $H$ as the retrieval half of the same
vector-lattice adjunction: any antitone Galois retrieval map must satisfy
\cref{eq:derivation-family}. It also suggests an observable diagnostic.
When $G_{\mathrm{img}}(u)\neq\varnothing$, define the witness discrepancy
\begin{equation}\label{eq:witness-discrepancy}
  \Delta(u)=1-\frac{|H(u)|}{|G_{\mathrm{img}}(u)|}.
\end{equation}
This measures the fraction of pooled matches lacking a same-site
witness. It is undefined for an empty pooled extent and is not a measure
of semantic error: dispersed evidence may be exactly what a query intends.

\subsection{When a vector suffices, and when a downset is needed}
\label{sub:downsets}
Write $\down v=\{u\in\U\mid u\leq v\}$. The set of queries witnessed
at some site of image $x$ is
\begin{equation}\label{eq:site-downset}
  s_{x}=\bigcup_{p\in P_{x}}\down z_{x,p}.
\end{equation}
It is a downset: whenever it contains a query, it contains all weaker
queries. Max pooling replaces it by the generally larger principal
downset $\down u_{x}$.

\begin{proposition}[Single-vector criterion for spatial evidence]
  \label{prop:principal-criterion}
  For a finite nonempty $P_{x}$, the downset $s_{x}$ is principal if and
  only if one site code dominates all the other site codes. In that case
  it equals $\down u_{x}$.
\end{proposition}
\begin{proof}
  If one site code dominates all the others, it is their coordinatewise
  maximum, and the union in \cref{eq:site-downset} is its principal downset.
  Conversely, suppose $s_{x}=\down v$. Every site code lies below $v$,
  so $u_{x}\leq v$. Because $v\in s_{x}$, some site $p$ satisfies
  $v\leq z_{x,p}\leq u_{x}$. Hence $v=z_{x,p}=u_{x}$, and that site code
  dominates all the others.
\end{proof}

The criterion describes a restrictive special case, not a generic
property of sparse representations. To retain same-site evidence, let
$\mathcal{D}(\U)$ be the set of all downsets of $\U$, including the
empty set, ordered by inclusion. Arbitrary unions and intersections are
its joins and meets; it is a complete lattice with bottom $\varnothing$
and top $\U$. Use $s_{x}$ as the description of image $x$ and define
\begin{equation}\label{eq:downset-derivations}
  F_{\mathrm{sp}}(A)=\bigcap_{x\in A}s_{x},\qquad
  G_{\mathrm{sp}}(u)=\{x\mid u\subseteq s_{x}\},
  \quad u\in\mathcal{D}(\U).
\end{equation}
Here $u$ denotes a downset-valued attribute, rather than a vector.

\begin{theorem}[Downset context for spatial evidence]
  \label{thm:spatial-context}
  The maps in \cref{eq:downset-derivations} form an antitone Galois
  connection. For every vector $v\in\U$,
  \begin{equation}\label{eq:principal-retrieval}
    G_{\mathrm{sp}}(\down v)=H(v).
  \end{equation}
  Thus same-site vector queries admit exact concept closure after passing
  to downset descriptions.
\end{theorem}
\begin{proof}
  The equivalence $A\subseteq G_{\mathrm{sp}}(u)$ if and only if
  $u\subseteq\bigcap_{x\in A}s_{x}$ proves the adjunction, just as in
  \cref{prop:adjunction}. Since $s_{x}$ is downward closed,
  $\down v\subseteq s_{x}$ is equivalent to $v\in s_{x}$, which is exactly
  the existence of a site code dominating $v$.
\end{proof}

This uses the established pattern-structure construction with richer
descriptions. It does not make same-site retrieval preserve vector joins:
$\down v\cup\down w$ generally differs from $\down(v\vee w)$.
The former can be witnessed at different sites, whereas the latter
requires a site satisfying the joint vector. Closure may likewise turn
a principal downset into a nonprincipal one.

The representation is computationally concrete. A finite union of
principal downsets is determined by its maximal generating vectors.
Intersections can be computed using
\begin{equation}\label{eq:downset-intersection}
  \left(\bigcup_{i}\down v_{i}\right)
  \cap\left(\bigcup_{j}\down w_{j}\right)
  =\bigcup_{i,j}\down(v_{i}\wedge w_{j}).
\end{equation}
Dominated generators can then be discarded. The number of undominated
generators can grow rapidly, so the pilot computes $H$ by direct site
comparisons and does not enumerate the downset concept lattice.
This preserves the distinction between an exact representation theorem
and a scalable mining algorithm.
